%==============================================================================
\section{Construction from public data}
\label{sec:construction}

\Cref{tab:sources} lists the data the kernel reads.
Every downloaded file is recorded with its address, retrieval date (9--12 November 2026) and SHA-256 digest.
Only identifiers and measured values are taken from each source; nothing is copied from a source's own inferences.

\begin{table}[t]
\centering
\caption{Data sources and what the kernel takes from each.}
\label{tab:sources}
\small
\begin{tabularx}{\linewidth}{@{}l>{\raggedright\arraybackslash}X@{}}
\toprule
Source & Content used \\
\midrule
Reactome \citep{gillespie2022reactome} & SBML export for \textit{Homo sapiens}, 2{,}825 pathway files; reactions, species, compartments, modifiers \\
ChEBI \citep{hastings2016chebi} & compounds, database accessions, ontology relations (conjugate acid/base, tautomer) \\
TECRDB \citep{goldberg2004tecrdb} & apparent equilibrium constants, reverse-transformed to the chemical reference state \\
eQuilibrator compound cache \citep{beber2022equilibrator} & microspecies (charge, protons, magnesium ions, relative energies) and masses per compound \\
SABIO-RK \citep{wittig2012sabio} & every human wild-type kinetic law: 5{,}682 $K_m$, 3{,}089 $k_{\mathrm{cat}}$ and 771 $K_i$ records \\
PaxDb \citep{wang2015paxdb,nagaraj2011hela} & whole-organism integrated abundances (3{,}048 proteins mapped); HeLa abundances \\
UniProt \citep{uniprot2023} & EC numbers, Pfam domains and families for 3{,}207 proteins; modified residues and active sites with evidence codes \\
AlphaFold Protein Structure Database \citep{varadi2022afdb} & predicted structures and per-residue confidence (pLDDT) \\
Protein Data Bank \citep{berman2000pdb} & experimental coordinates \\
PDBe best structures and SIFTS \citep{velankar2013sifts} & choice of experimental chain; residue numbering mapped to UniProt \\
Rhea \citep{bansal2022rhea} & mappings to Reactome, reaction directions, participants with coefficients \\
MetaNetX \citep{moretti2021metanetx} & chemical cross-references (ChEBI to KEGG) \\
Human-GEM \citep{robinson2020humangem} & cytosolic reactions, metabolite identifiers, reversibility \\
Human Protein Atlas \citep{thul2017hpa} & main and additional subcellular locations \\
COMPARTMENTS \citep{binder2014compartments} & cytosol (GO:0005829) confidence scores \\
M-CSA \citep{ribeiro2018mcsa} & catalytic residues in UniProt numbering \\
DisProt \citep{quaglia2022disprot} & experimentally observed disordered regions \\
\bottomrule
\end{tabularx}
\end{table}

\subsection{The cytosolic corpus}
\label{sec:corpus}

The Reactome export is streamed once.
A reaction is kept when every species it references, whether reactant, product or modifier, lies in the compartment named \emph{cytosol}.
Reactions are de-duplicated across pathway files, and the pathways each one appears in are recorded.
The corpus has 2{,}914 reactions over 5{,}344 species: 4{,}103 macromolecules, 796 small molecules and 445 other entities.

\subsection{Species identity}
\label{sec:identity}

Each small molecule is mapped from its ChEBI accession to a KEGG compound \citep{kanehisa2000kegg}, because TECRDB and eQuilibrator are keyed by KEGG.
The basis of each mapping is recorded.
It is \emph{stated} when the ChEBI record itself carries exactly one KEGG accession (239 species).
It is \emph{supplied} under convention \conv{conjugate-closure} when the accession is found through the conjugate-acid, conjugate-base or tautomer relations (309).
It is \emph{absent} when neither the record nor its closure carries one (235), and \emph{underdetermined} when several candidates remain (13).
Protonation states share a reaction-level KEGG identifier, and TECRDB is protonation-agnostic, which justifies the closure.
An underdetermined species keeps all candidates, and a reaction match is accepted only if exactly one combination of candidates matches a TECRDB reaction.
The proton (C00080) is dropped on both sides of every reaction: the chemical reference state gives it zero formation energy, and Reactome states it inconsistently.

\subsection{Reaction thermodynamics}
\label{sec:thermo}

\paragraph{Pooling.}
TECRDB rows are kept when their temperature lies in 293.15--313.15\,K and their quality grade is A, B or C.
Each row's chemical-reference $\DG$ is put into an orientation-free canonical form (a key and a sign).
Rows from one study are averaged, and the reaction's value is the mean of its study means.
Replicate variability is pooled over reactions measured by two or more studies,
\begin{equation}
  \sigma_{\mathrm{pool}}^2 = \frac{\sum_{\rho}\sum_{i\in\rho}(g_{\rho i}-\bar g_\rho)^2}{\sum_\rho (n_\rho-1)},
  \qquad \sigma_\rho = \sigma_{\mathrm{pool}}/\sqrt{n_\rho},
\end{equation}
which gives $\sigma_{\mathrm{pool}}=2.68\kJ$ over 418 reactions.
Tolerances use $q=2.576$ (two-sided Gaussian at $\alpha=0.01$) and, as a distribution-free alternative, the Chebyshev value $q=\alpha^{-1/2}=10$.

\paragraph{Matching.}
A Reactome reaction matches a TECRDB reaction when their canonical keys agree.
A match that holds only once water is balanced is accepted under convention \conv{water-balance}.
A reaction with a macromolecular participant is \vd{refused}, because TECRDB states small-molecule reactions only.

\paragraph{Three-valued cycle test.}
Let the matched reactions be the columns of a small-molecule stoichiometric matrix $S$.
An integer basis of $\ker S$ is computed by exact rational elimination.
For a basis vector $z$, the cycle sum $H=\sum_j z_j g_j$ must vanish.
Its rounding tolerance is $\varepsilon_{\mathrm{num}}=\gamma_{|z|}\sum_j|z_j g_j|$, with $\gamma_n=nu/(1-nu)$ and unit roundoff $u=2^{-53}$ \citep{higham2002accuracy}.
Its data tolerance is $\varepsilon_{\mathrm{data}}=q\,(\sum_j z_j^2\sigma_j^2)^{1/2}$.
The verdict is
\begin{equation}
  \text{consistent if } |H|\le\varepsilon_{\mathrm{num}};\quad
  \text{undecidable if } \varepsilon_{\mathrm{num}}<|H|\le\varepsilon_{\mathrm{num}}+\varepsilon_{\mathrm{data}};\quad
  \text{inconsistent otherwise.}
\end{equation}
A planted defect in a synthetic cycle is flagged, and a consistent cycle passes (\cref{app:tests}).

\paragraph{Hess's law.}
Let $N$ be the compounds-by-reactions matrix of all measured TECRDB reactions, with values $g$ and standard errors $s$.
Formation energies are fixed by the measurements only up to the left null space of $N$.
The weighted least-squares estimate is
\begin{equation}
  \hat\mu = \arg\min_\mu \bigl\| W^{1/2}(N^\top\mu-g)\bigr\|_2,\qquad W=\operatorname{diag}(1/s^2).
\end{equation}

\begin{proposition}[Determinacy of Hess's law]
\label{prop:hess}
For a target reaction $b$, the value $b^\top\hat\mu$ is the same for every minimiser $\hat\mu$ if and only if $b\in\im N$.
Its variance is $b^\top (NWN^\top)^{+} b$.
\end{proposition}
\begin{proof}
Two minimisers differ by $u$ with $N^\top u=0$.
If $b=Ny$, then $b^\top u=y^\top N^\top u=0$.
If $b\notin\im N$, then $b$ has a component in $\ker N^\top$, and adding a multiple of that component to $\hat\mu$ leaves the residual unchanged while changing $b^\top\hat\mu$.
The variance follows from the weighted normal equations.
\end{proof}

A target in the span is \vd{answered} under convention \conv{hess-span}; a target with a compound that was never measured, or outside the span, is \vd{absent}, and its blocker says which.
Cycle inconsistency among the measured reactions shows up as the fit's $\chi^2$, which is reported in \cref{sec:res-thermo}.

\begin{proposition}[Falsifiability guard]
\label{prop:falsifiable}
If every reaction energy in a cycle is computed from species energies, $g=N^\top f$ for some $f$, then every cycle sum vanishes identically.
\end{proposition}
\begin{proof}
For $z\in\ker N$, $z^\top g = z^\top N^\top f = (Nz)^\top f = 0$.
\end{proof}

A cycle test built only from derived or estimated energies therefore cannot fail, and the kernel refuses it with the blocker ``non-falsifiable''.
A cycle with at least one stated or supplied measurement is judged as above.
This guard is what separates the kernel's cycle test from a check on component-contribution estimates \citep{noor2013component}, which are consistent by construction.

\subsection{The medium: Legendre transform over the cytosolic window}
\label{sec:legendre}

The medium enters only through the Legendre transform of each compound \citep{alberty1998calculation}.
Let compound $i$ have microspecies $j$ of charge $z_{ij}$, with $n^{\mathrm{H}}_{ij}$ protons and $n^{\mathrm{Mg}}_{ij}$ magnesium ions.
Let $\Delta\Delta G_{ij}$ be its energy relative to the compound's reference microspecies.
The transform is then
\begin{align}
  t_i &= -RT\ln\sum_j \exp\Bigl(-\frac{\Delta\Delta G_{ij}+\tau_{ij}}{RT}\Bigr),\\
  \tau_{ij} &= n^{\mathrm{H}}_{ij}\,RT\ln10\,\mathrm{pH} + n^{\mathrm{Mg}}_{ij}\bigl(RT\ln10\,\mathrm{pMg}-\Delta_f G_{\mathrm{Mg}}(T)\bigr)\nonumber\\
  &\quad - \alpha(T)\frac{\sqrt I}{1+1.6\sqrt I}\bigl(z_{ij}^2-n^{\mathrm{H}}_{ij}-4n^{\mathrm{Mg}}_{ij}\bigr),
\end{align}
a function of pH, $I$, $T$ and pMg, with the extended Debye--H\"uckel coefficient $\alpha(T)$ cubic in $T$, as in eQuilibrator \citep{beber2022equilibrator}.
A reaction's transformed value is $\DGp=\DG+\sum_i\nu_i t_i$.
Our implementation was checked against eQuilibrator's own transforms row by row (\cref{sec:res-medium}).

The cytosolic window (convention \conv{medium-window}) is pH 6.8--7.6, ionic strength 0.1--0.25\,M and pMg 2.8--3.5, at 298.15\,K, sampled on a grid of $5\times3\times3=45$ conditions.
The temperature is 298.15\,K because the TECRDB values are not temperature-corrected.
The concentration window is 1\,\textmu M--10\,mM for every metabolite (convention \conv{concentration-window}).
Crowding enters as an accessible-volume fraction $a\in[0.7,1]$ with $\mu=\mu^\circ+RT\ln(c/a)$ for every small solute (convention \conv{accessible-volume}).
It shifts a reaction's affinity by $-RT\ln a\sum_i\nu_i$, water excluded, so only reactions that change the number of solute particles feel it.

\subsection{Direction verdicts}
\label{sec:direction}

For a reaction with a chemical-reference value $\DG\pm\sigma$, let $[D_-,D_+]$ be the range of $\DGp$ over the medium grid widened by $q\sigma$.
Let $\ln Q_\pm$ be the extreme reaction quotients over the concentration window (products at the top of their windows and reactants at the bottom, or the reverse).
The verdict is
\begin{equation}
  \text{forward if } D_+ + RT\ln Q_+ < 0;\qquad
  \text{reverse if } D_- + RT\ln Q_- > 0;\qquad
  \text{\vd{undecidable} otherwise.}
\end{equation}
An undecidable verdict names its blocker.
The blocker is \emph{concentration} when the sign of $\DGm$, the value with every participant at 1\,mM, is fixed across the medium window, so only concentrations leave the sign open.
It is \emph{medium} when the medium window alone admits both signs.
The base stage, before the medium transform existed, used TECRDB's transformed study means in the pH window, without $q\sigma$ and without the ionic-strength and magnesium ranges.

\subsection{Carrier and slot: thermodynamics for protein-modifying reactions}
\label{sec:slot}

TECRDB refuses every reaction with a protein participant.
Many such reactions are phosphorylations in which the same protein entity appears on both sides and only small molecules and modified residues change.
We decompose such a reaction into a carrier, the polypeptide whose identity is inert, and slots, the side chains that change:
\begin{align}
  \text{kinase:}\quad & \mathrm{X}+\mathrm{ATP}\to\mathrm{pX}+\mathrm{ADP}\qquad(\mathrm X\in\{\mathrm{Ser},\mathrm{Thr},\mathrm{Tyr}\}),\\
  \text{phosphatase:}\quad & \mathrm{pX}+\mathrm{H_2O}\to\mathrm X+\mathrm{P_i}.
\end{align}
The slot reaction is small-molecule chemistry and can be put to TECRDB again, through Hess's law.
Four conventions are named on every such value:
\conv{carrier-slot} (the carrier contributes nothing to $\DG$);
\conv{free-residue} (a residue in a protein reacts like the free amino acid);
\conv{slot-from-name} (new sites are read from Reactome species names); and
\conv{hess-span}.
A reaction whose names state a number of changed sites inconsistent with its kinase or phosphatase steps is \vd{underdetermined}, and the blocker states both counts.

\subsection{Catalysts as closed ladders}
\label{sec:enzyme}

In the base graph a catalyst is a modifier, its node law reads $0=0$ and its state is underdetermined.
Under convention \conv{enzyme-cycle-2}, every catalysed junction $r$ with catalyst $E$ is expanded into a binding rung and a catalytic rung:
\begin{equation}
  r^{\mathrm{bind}}:\ E+\text{reactants}\to E{\cdot}r,\qquad
  r^{\mathrm{cat}}:\ E{\cdot}r\to E+\text{products}.
\end{equation}

\begin{proposition}[Enzyme totals]
\label{prop:enzyme}
The expansion leaves every net stoichiometry unchanged.
If $E$ takes part in no reaction other than as a catalyst, $E+\sum_r E{\cdot}r$ is a left null vector of the expanded matrix, i.e.\ a conservation law.
\end{proposition}
\begin{proof}
Summing the two rungs cancels $E$ and $E{\cdot}r$ and restores the net reaction.
In $r^{\mathrm{bind}}$ the total loses one $E$ and gains one $E{\cdot}r$; in $r^{\mathrm{cat}}$ the reverse.
No other column contains $E$ or $E{\cdot}r$, so every column sums to zero against the total.
\end{proof}

When $E$ is also made, consumed or modified by another reaction, the total is broken, and the kernel classifies what breaks it (\cref{sec:res-extend}).

\subsection{Proteins as contact graphs}
\label{sec:proteins}
\label{sec:attach}

Each protein with a slot is represented by its AlphaFold model as a residue graph.
The weight between two residues is the number of heavy-atom pairs closer than 4.5\,\AA.
A residue's medium weight is the number of its heavy atoms with more than 1\,\AA$^2$ of solvent-accessible surface, computed by the Shrake--Rupley method with a 1.4\,\AA\ probe and 192 points per atom \citep{shrake1973environment}.
A buried residue still receives medium weight 1 (convention \conv{medium-floor-residue}), so every residue touches the medium.
The vectorised surface computation was checked against the brute-force loop on 300 sampled atoms and against the analytic area of an isolated sphere.

A protein is attached through its slots.
Each named, positioned site of a carrier reaction is located in the protein's UniProt sequence, the residue type is checked, and an edge of weight 1 joins the junction to the residue vertex.
By \cref{prop:monotone}, this expansion cannot lower any existing item's separation cost.

\subsection{Kinetics: retention and locality}
\label{sec:kinetics}

We treat the cytosol as a chromatographic column in which free diffusion is transit and a stay bound to a protein is a contact.
A molecule that alternates between a free and a bound state moves, on long scales, at its free velocity divided by $1+k$, with retention factor $k$ \citep{giddings1991unified}.
In the circuit, fast binding raises the capacitance in \cref{eq:node} from $c/RT$ to $(1+k)c/RT$.
For a metabolite $S$ bound by enzymes $E$,
\begin{equation}
  k_S=\sum_E\frac{[E]_{\mathrm{free}}}{K_d(E,S)}\ \ge\ \sum_E\frac{[E]_{\mathrm{free}}}{K_m(E,S)},
\end{equation}
because $K_m=(k_{\mathrm{off}}+k_{\mathrm{cat}})/k_{\mathrm{on}}\ge K_d$.
Since $[E]_{\mathrm{free}}\le[E]_{\mathrm{total}}$, the $K_m$-based value is reported per unit free enzyme and, under convention \conv{protein-ppm-to-molar}, at total enzyme.
That convention converts PaxDb parts per million to molar using 2--6\,mM total protein molecules \citep{milo2013protein}.

The carrier can be dropped from the description only when contacts are local, i.e.\ when one item's binding does not depend on how many others are present.
An enzyme's sites are shared by all its substrates, so the readout for one substrate departs from its dilute limit by the factor $1/(1+\sum_S c_S/K_S)$.
We report, for each enzyme, the concentration $c^\ast=0.1/\sum_S K_S^{-1}$ at which that departure reaches 10\% with every substrate at $c^\ast$.

\subsection{Completion}
\label{sec:complete}

A completion of source $S$ into target $P$ is a set of junctions, each fired in an allowed direction, that turns $S$ into $P$.
The other inputs are the medium and the 2{,}241 proteins that no reaction makes, taken as present (convention \conv{scope-from-source}); a complex is made, so a complex holding a small ligand is not a free source of that ligand.
The search computes the forward scope, the closure of the source under junctions whose inputs are all available.
It then works backward from the target, keeping only junctions in the scope that produce a needed species and requiring their inputs in turn.
A blocked species is reached but never spent as an input, so a completion of a target cannot pass through the target itself.
The outcomes map onto verdicts: no completion is \vd{certified-empty} relative to the allowed directions; exactly one is \vd{answered} (an order-determined ladder); several are \vd{underdetermined}, and the branch points are named.

Three direction semantics are compared:
\emph{as annotated} (Reactome's direction only);
\emph{reversible} (both directions unless thermodynamics decides); and
\emph{fixed at 1\,mM} (as reversible, but also fixing the sign wherever $\DGm$ has a fixed sign across the medium window).
The third is a counterfactual: it shows what measured concentrations near 1\,mM would add (\cref{sec:res-complete}).

\subsection{Modifiers, flux dividers and the redox ladder}
\label{sec:modifiers}

\paragraph{Modifiers.}
An inhibitor edge scales the junction's prefactor by $1/(1+[I]/K_i)$ over the concentration window, as an additive term in $\ln k_r$.
Without a $K_i$ in SABIO-RK the stack is \vd{absent}; a macromolecular binding partner is \vd{refused}, because no $K_i$ applies to it.
Reactome's loss-of-function variant reactions are attached as modifiers on their wild-type edge, not as competing siblings (convention \conv{loss-of-function-variant}).

\paragraph{Flux dividers.}
At a small-molecule branch point, branch $i$'s share in the dilute limit is
\begin{equation}
  \phi_i=\frac{(k_{\mathrm{cat}}/K_m)_i[E_i]}{\sum_j (k_{\mathrm{cat}}/K_m)_j[E_j]}
\end{equation}
(convention \conv{dilute-limit-share}).
The verdict is \vd{answered} only when every sibling has $k_{\mathrm{cat}}$, $K_m$ and abundance; otherwise the missing sibling is named.
SABIO-RK records name no substrate on $k_{\mathrm{cat}}$, so each $k_{\mathrm{cat}}$ is paired with the $K_m$ values of the same kinetic law.

\paragraph{Redox ladder.}
Couple potential differences $\Delta E'^\circ=-\DGp/(nF)$ come from the Hess span with the Legendre transform at pH 7, $I=0.25$\,M and 298.15\,K, without magnesium (convention \conv{textbook-standard-conditions}).
For intramolecular electron transfer, the Moser--Dutton ruler gives \citep{moser1992nature,page1999natural}
\begin{equation}
  \log_{10}k_{\mathrm{ET}} = 13-(1.2-0.8\rho)(R-3.6)-3.1\frac{(\Delta G+\lambda)^2}{\lambda},
\end{equation}
with edge distance $R$ in \AA\ (taken as 3.6\,\AA\ when shorter), packing density $\rho=0.76$, reorganisation energy $\lambda=0.7$\,eV and driving force $\Delta G=0$ (convention \conv{moser-dutton-defaults}).

\subsection{Labels and stoichiometric signatures}
\label{sec:labels}

Junctions receive the EC class of their catalyst, truncated to levels 1--4, along with Pfam and gene labels from UniProt.
Labels are compared with the classes of colour refinement against 200 shuffles of the same labels among the same junctions.
Independently, each reaction receives a class from its stoichiometry alone by six fixed rules.
A reaction is an \emph{oxidoreductase} when an NAD, NADP, flavin adenine dinucleotide, glutathione or thioredoxin couple changes state.
It is a \emph{hydrolase} when water is consumed and there are more small products than reactants other than water.
It is a \emph{ligase} when a nucleoside triphosphate is split alongside a bond formation.
It is a \emph{transferase} for ATP to ADP with a phosphorylated product, for S-adenosylmethionine to S-adenosylhomocysteine, or for a coenzyme~A ester.
It is an \emph{isomerase} when one small molecule goes in and one with the identical formula comes out.
It is a \emph{lyase} when one small molecule goes in, two come out and no water is involved.
When several rules match, the verdict is \vd{underdetermined}; a reaction with a macromolecular participant is \vd{refused}.

\subsection{Cross-checks against curated databases}
\label{sec:xref}

\paragraph{Rhea.}
Rhea's mapping to Reactome gives, for each kernel reaction, the directional Rhea reaction (left-to-right or right-to-left) that matches Reactome's orientation.
Participants with coefficients are read from Rhea's public graph-query service, where the coefficient is encoded in the predicate.
The curated direction is then compared with the sign of $\DGm$ wherever that sign is fixed across the medium window.

\paragraph{MetaNetX.}
For a species left without a KEGG compound by \conv{conjugate-closure}, a KEGG compound reached through a MetaNetX identifier is supplied under convention \conv{mnx-xref}.

\paragraph{Human-GEM.}
Kernel reactions whose small-molecule participants all carry KEGG identifiers are matched to Human-GEM cytosolic reactions by participant set, ignoring the proton.
Coefficients and reversibility are compared, the coefficients up to an overall sign.

\paragraph{Localisation.}
A catalyst has cytosol support when the Human Protein Atlas lists the cytosol as a main or additional location, or COMPARTMENTS gives it a cytosol confidence of at least 3.

\paragraph{Residue evidence.}
Located slots are compared with UniProt modified-residue features of phospho type, and the Evidence and Conclusion Ontology (ECO) code of each matched feature is read \citep{giglio2019eco}.
For each protein with residue mismatches, the single numbering offset in $[-120,120]$ that moves the most mismatched sites onto a phospho-residue of the named type is reported.
M-CSA catalytic residues are compared with residues of the same amino-acid types in the same AlphaFold contact graphs.
DisProt regions split the phosphosite comparison of \cref{sec:res-folding} into inside and outside measured disorder.

\paragraph{Experimental structures.}
For high-confidence AlphaFold models, the first X-ray chain covering at least 80\% of the model is taken from PDBe's best-structure list, and its numbering is mapped to UniProt through the Structure Integration with Function, Taxonomy and Sequences resource (SIFTS).
Contact sets (residue pairs with $|i-j|\ge3$ in contact) are compared by the Jaccard index, against a control that shifts one numbering by three residues.

\subsection{Loaders for measurements not yet present}
\label{sec:loaders}

Two further sources are wired as loaders.
A per-compound concentration window can be declared in a table with columns \texttt{kegg}, \texttt{low\_M}, \texttt{high\_M}, \texttt{cell\_type} and \texttt{source}; it replaces the global window for the compounds it lists (convention \conv{measured-window}).
The BRENDA release can be read from its JSON export \citep{chang2021brenda}, with a $k_{\mathrm{cat}}$ paired to a $K_m$ only when both come from the same protein, substrate and reference (convention \conv{brenda-same-reference}).
Without the files, each loader returns \vd{absent} and names the file it needs.
At the time of writing, neither file was present (\cref{sec:limitations}).
